\chapter{Runtime internals: a guided tour}
\label{ch:runtime}

\begin{aside}
We have described \code{maya::run} as ``the runtime''. Plural,
really --- there are several cooperating subsystems behind the
single \code{run} call. This chapter is a guided tour of the
internals: what each piece does, where it lives in the source,
and how they fit together. By the end you should be able to
open any header in \file{maya/include/maya/app/} and know what
you are reading.
\end{aside}

\section{The roadmap}

A \maya{} runtime is six cooperating pieces:

\begin{enumerate}[leftmargin=*]
  \item \textbf{Context.} The bundle of platform handles,
    config, capabilities, and stats.
  \item \textbf{Event loop.} The driver that waits for events,
    dispatches them, and runs the pipeline.
  \item \textbf{Event source.} The platform abstraction (epoll,
    kqueue, WaitForMultipleObjects) that produces events.
  \item \textbf{Subscription manager.} The diff of declared
    subscriptions versus active OS-level wires.
  \item \textbf{Command interpreter.} The interpreter that walks
    \code{Cmd<Msg>} variants and effects them.
  \item \textbf{Renderer.} The pipeline that turns elements into
    ANSI bytes.
\end{enumerate}

Each lives behind its own header. \code{maya::run} composes them.

\section{Context}

\file{maya/include/maya/app/context.hpp} defines:

\begin{lstlisting}
struct Context {
    Config            config;
    platform::Terminal& terminal;
    TerminalCaps      caps;
    SignalDispatcher& signals;
    WakeFd&           wake;
    Stats             stats;
    std::optional<Logger> logger;
};
\end{lstlisting}

The context is constructed at startup and passed by reference
into every subsystem. Subsystems read what they need; the
context owns nothing they own.

\subsection*{Why a context, not globals?}

Globals would work, but:

\begin{itemize}[leftmargin=*]
  \item Two simultaneous \maya{} runtimes in one process (a test
    runner, for example) would clobber each other.
  \item Threading is harder to reason about with hidden state.
  \item Tests cannot vary the context without resetting globals.
\end{itemize}

A context is a small ceremony for a real benefit.

\section{Event loop}

The driver in \file{maya/include/maya/app/app.hpp}:

\begin{lstlisting}
template <Program P>
int run(Config cfg) {
    auto ctx = make_context(cfg);
    auto loop = EventLoop{ctx};
    auto renderer = Renderer{ctx};
    auto interpreter = Interpreter{ctx, loop};

    auto [model, init_cmd] = P::init();
    interpreter.interpret(init_cmd);

    while (loop.running()) {
        renderer.frame(P::view(model));
        loop.update_subscriptions(P::subscribe(model));

        Event ev = loop.wait();
        if (auto msg = to_message<typename P::Msg>(ev)) {
            try {
                auto [m2, c2] = P::update(model, *msg);
                model = std::move(m2);
                interpreter.interpret(c2);
            } catch (...) {
                log_update_exception(ctx, std::current_exception());
            }
        }
    }
    return 0;
}
\end{lstlisting}

The pattern is the loop we walked in Chapter~\ref{ch:elm}, now
with the supporting classes named.

\subsection*{Composition over inheritance}

\code{EventLoop}, \code{Renderer}, and \code{Interpreter} are
plain classes that take a \code{Context\&} in their constructor.
There is no \code{class App} hierarchy. The runtime is the
\code{run} function plus the helpers it instantiates.

\begin{idea}
Composition keeps the runtime debuggable. Each subsystem is a
file; each file is a class; each class has a clear contract.
You can read the runtime by reading six files in order, without
following virtual dispatches.
\end{idea}

\section{Event source}

The OS-level pump (Chapter~\ref{ch:platform}). It produces
\code{Event} values: \code{KeyPress}, \code{MouseEvent},
\code{Resize}, \code{Wakeup}, \code{Timeout}.

The event loop calls \code{source.wait(timeout)} and translates
the result into a message via the active subscription set.

\subsection*{The Event variant}

\begin{lstlisting}
using Event = std::variant<
    KeyPress, MouseEvent, Resize, Paste,
    Timeout, Wakeup, Unknown>;
\end{lstlisting}

The same overload pattern from Chapter~\ref{ch:elements} matches
each kind. The subscription manager (next section) handles the
match.

\section{Subscription manager}

The subscriber tracks which \code{Sub} alternatives are active
and how to translate events into messages:

\begin{lstlisting}
class SubscriptionManager {
public:
    void update(const Sub<Msg>& new_sub);
    std::optional<Msg> dispatch(const Event& ev);
};
\end{lstlisting}

\code{update} takes the new \code{Sub} from
\code{P::subscribe(model)} and diffs it against the previous one.
Subscriptions added are installed; subscriptions removed are
torn down. Subscriptions present in both are kept.

\code{dispatch} takes an event and returns a message if any
subscription matches. The matching is by event kind and any
subscription-specific filters.

\subsection*{Diffing subscriptions}

The diff uses structural equality on a flattened list of
(kind, parameters) tuples. Two \code{Every(50ms, Tick\{\})}
subscriptions are equal regardless of how they were constructed.
The fingerprint comparison is O(n) in subscription count, but
subscription counts are tiny (typically under ten).

\section{Command interpreter}

\file{maya/include/maya/app/} also houses the interpreter:

\begin{lstlisting}
class Interpreter {
public:
    void interpret(const Cmd<Msg>& cmd);

private:
    void interpret_none(None);
    void interpret_quit(Quit);
    void interpret_after(After<Msg>);
    void interpret_task(Task<Msg>);
    void interpret_batch(Batch<Msg>);
    // ... one per alternative
};
\end{lstlisting}

Each \code{interpret\_X} performs the real-world side effect
corresponding to that alternative. \code{interpret\_quit} sets a
flag on the event loop. \code{interpret\_after} schedules a
timer on the event source. \code{interpret\_task} hands the
callable to the worker pool.

The dispatch is a single \code{std::visit}:

\begin{lstlisting}
void Interpreter::interpret(const Cmd<Msg>& cmd) {
    std::visit([this](auto& alt) { this->do_interpret(alt); }, cmd.v);
}
\end{lstlisting}

Adding a new \code{Cmd} alternative is two steps: add to the
variant, add an \code{interpret} overload. No existing callers
change.

\subsection*{The worker pool}

\code{Task} and \code{IsolatedTask} run on threads. \maya{}
ships with a small fixed-size pool plus the option for isolated
threads. The pool size is configurable (default: thread::hardware
concurrency divided by two; minimum one).

A pool worker, when its callable returns, writes the resulting
message to a lock-free queue and signals the wake fd. The main
loop's event source picks up the wake; \code{dispatch} pulls the
message from the queue.

\begin{aside}
The wake-fd pattern is what makes \maya{} able to deliver
worker-thread messages without polling. The cost is one byte
written and one byte read per cross-thread message.
\end{aside}

\section{The renderer}

\file{maya/include/maya/render/renderer.hpp} houses:

\begin{lstlisting}
class Renderer {
public:
    explicit Renderer(Context&);
    void frame(Element root);
    void force_redraw();
    void set_caps(TerminalCaps);
};
\end{lstlisting}

\code{frame(Element)} runs the six-stage pipeline from
Chapter~\ref{ch:pipeline}: layout, paint, diff, serialise,
write. \code{force\_redraw} discards the cache. \code{set\_caps}
updates capability detection (e.g. on terminal-type change).

The renderer owns the previous canvas (for the diff) and the
output buffer. Each frame is amortised; the buffer's capacity
persists across frames.

\section{Shutdown}

When \code{Cmd::Quit} is interpreted, the event loop sets its
running flag to false. The main \code{while} exits on the next
iteration. Destructors restore the terminal:

\begin{enumerate}[leftmargin=*]
  \item \code{Renderer}'s destructor flushes any pending bytes.
  \item \code{Interpreter}'s destructor joins worker threads.
  \item \code{EventLoop}'s destructor closes the event source.
  \item \code{Terminal}'s destructor calls \code{leave\_raw\_mode}.
\end{enumerate}

The order matters: flush bytes \emph{before} restoring raw mode
(so the final frame appears), join workers \emph{before}
destroying the queue (so writes do not race), and so on. RAII
takes care of this if the construction order is right.

\subsection*{Signal-driven shutdown}

If the user presses Ctrl-C, the signal handler writes a special
\code{InterruptedSignal} message to the queue. \code{update}
sees it (most apps handle it as ``quit''); the rest is the
normal shutdown path.

If the signal handler cannot reach \code{update} (e.g. the
runtime is stuck in a long computation), the OS will eventually
kill the process. The terminal restore happens in
\code{atexit}-registered handlers as a last resort.

\section{Concurrency model summary}

\begin{itemize}[leftmargin=*]
  \item Main thread: event loop, renderer, interpreter,
    \code{update}, \code{view}, \code{subscribe}.
  \item Worker threads: \code{Task} callables.
  \item Isolated threads: \code{IsolatedTask} callables.
  \item Signal-disposition thread (on Win32): the console
    control handler.
\end{itemize}

The handoff between worker and main is via the wake fd and the
message queue. The main thread is the only one that touches
\code{Model}, \code{view}, or the renderer.

\begin{warning}
Do not share mutable state across the worker/main boundary
without synchronisation. The whole architecture is built on the
idea that \code{Model} is owned by the main thread; if you
violate that, races are likely.
\end{warning}

\section{Reading the source}

A suggested order for reading \maya{}'s runtime:

\begin{enumerate}[leftmargin=*]
  \item \file{maya/include/maya/app/app.hpp} (the \code{run}
    function and the main loop).
  \item \file{maya/include/maya/app/context.hpp}.
  \item \file{maya/include/maya/app/sub.hpp}.
  \item \file{maya/include/maya/core/cmd.hpp}.
  \item \file{maya/include/maya/app/events.hpp}.
  \item \file{maya/include/maya/render/renderer.hpp}.
\end{enumerate}

Six files. Combined, around 2000 lines of (heavily commented)
source. That is the entire runtime.

\section*{Workshop}

\begin{enumerate}[leftmargin=*]
  \item Open \file{maya/include/maya/app/app.hpp} and identify
    every subsystem class instantiated in \code{run}.
  \item Trace what happens when the user presses Ctrl-C from
    the terminal to the program's exit. Note each thread
    boundary crossed.
  \item Try implementing a simple\code{trace\_all} mode that
    logs every event, message, and command. Decide which
    subsystem you would hook into.
\end{enumerate}

\begin{recap}
The runtime is a composition of six subsystems: context, event
loop, event source, subscription manager, command interpreter,
renderer. Each lives behind its own header; \code{run} composes
them. The main thread owns the model; worker threads communicate
via a wake-fd and message queue. RAII handles shutdown order.
The whole thing is around 2000 lines of source.
\end{recap}
